\section{Introduction}
\label{sec:intro}

An equation over a group permits coefficients from that group and asks
whether the unknowns can be realized in a larger group without identifying
distinct coefficients. The exponent sums record the dependence on the
unknowns after abelianization. A system whose exponent rows are independent
is called nonsingular. We prove that every such
finite system is solvable over every coefficient group.

Let $G$ be a group and let $F_n=F(x_1,\ldots,x_n)$ be the free group on
$n$ generators. A system over $G$ consists of words
$w_1,\ldots,w_m\in G*F_n$, interpreted as equations $w_i=1$. A solution
over $G$ is a group $H$ containing an isomorphic copy of $G$, together
with elements $h_1,\ldots,h_n\in H$ at which all the words evaluate
to the identity. Define the \emph{exponent-sum matrix}
$A=(a_{ij})\in M_{m\times n}(\mathbb Z)$ by the homomorphism
\[
 G*F_n\longrightarrow\mathbb Z^n,
 \qquad G\longmapsto0,\quad x_j\longmapsto e_j:
 \qquad w_i\longmapsto(a_{i1},\ldots,a_{in}).
\]
The system is \emph{nonsingular} if $\operatorname{rank}_{\mathbb Q}A=m$.
Here $\langle\!\langle w_1,\ldots,w_m\rangle\!\rangle$ denotes normal
closure in $G*F_n$.

\begin{theorem}[Nonsingular systems]\label{thm:main}
Let $G$ be any group, let $1\le m\le n$ be integers, and let
$w_1,\ldots,w_m\in G*F_n$. If their exponent-sum matrix has row rank
$m$ over $\mathbb Q$, then the canonical homomorphism
\[
 G\longrightarrow
 (G*F_n)/\langle\!\langle w_1,\ldots,w_m\rangle\!\rangle
\]
is injective. Equivalently, the system has a simultaneous solution in
an overgroup of $G$.
\end{theorem}

The equivalence follows from the universal property of the displayed
quotient. If the canonical map is injective, the quotient itself supplies
the overgroup and the images of the free generators supply the solution.
Conversely, any solution induces a homomorphism from the quotient whose
restriction to $G$ is injective. Thus solving the equations amounts to
showing that their normal closure kills no nonidentity element of $G$.
Theorem~\ref{thm:main} imposes no countability, finiteness, or torsion
condition on $G$. The same conclusion for arbitrary sets of equations
and variables with independent exponent rows follows by a finite-relation
argument in Remark~\ref{rem:infinite-systems}.

\subsection{History and significance}
The general solvability assertion for nonsingular systems is commonly
called Howie's conjecture; see \cite[Introduction]{KlyachkoMikheenkoRomankov2024}
for this formulation and terminology. Its one-equation, one-variable
case is the Kervaire--Laudenbach conjecture: a nonzero exponent sum
should guarantee coefficient injectivity. This is stronger than the
nontriviality assertion usually called the Kervaire conjecture, which
asks whether adjoining one generator and imposing one relation can
annihilate a nontrivial group. Chen explains these formulations and their
connection with high-dimensional knot groups
\cite[Section~1]{Chen2026}. Theorem~\ref{thm:main} establishes the
finite nonsingular-system assertion for arbitrary coefficient groups.

The foundational system theorem is due to Gerstenhaber and Rothaus
\cite{GerstenhaberRothaus1962}. For a square system over a compact
connected Lie group, they compute the degree of the word map from the
determinant of its exponent-sum matrix. A nonzero determinant therefore
forces surjectivity and a simultaneous solution in that Lie group.
Their finite-group theorem also produces a finite solution overgroup,
using arithmetic specialization and reduction over finite fields.
A nonsingular rectangular
system reduces to the square case by retaining columns of a nonzero
maximal minor and setting the other variables equal to the identity. Their
argument established the usefulness of compact Lie groups, cohomology,
and degree theory in a problem stated entirely in group-theoretic terms.

A different line of work uses the topology of relative presentations.
Howie proved solvability of finite independent systems over locally
indicable groups \cite[Corollary~4.2]{Howie1981}; locally indicable
means that every nontrivial finitely generated subgroup maps onto
$\mathbb Z$. Klyachko's theorem treats a single equation in one variable
over a torsion-free group when the exponent sum is $\pm1$
\cite{Klyachko1993}. This unimodular hypothesis is stronger than a
nonzero exponent sum. Chen later recovered that theorem by estimating
the complexity of surfaces in HNN extensions
\cite[Theorem~6.9]{Chen2026}. These results illustrate two complementary
sources of control: algebraic conditions on the coefficient group and
topological restrictions on diagrams witnessing a kernel element.

The compact-unitary method also passes to metric ultraproducts. Pestov
observed its application to hyperlinear groups
\cite[Corollary~10.4]{Pestov2008}, and Nitsche and Thom give the
nonsingular-system statement explicitly
\cite[Theorem~1.2 and Lemma~2.1]{NitscheThom2022}. Here a hyperlinear,
or Connes-embeddable, group is one that embeds into a metric ultraproduct
of unitary groups with their normalized Hilbert--Schmidt metrics.
Their stronger Theorem~1.3 replaces nonsingularity by vanishing second
homology of a covering of the presentation complex obtained after
deleting the coefficients. For the presentation complex itself,
vanishing second homology is precisely independence of the exponent
rows. Their covering criterion therefore reaches beyond nonsingular
systems while retaining the hypothesis on the coefficient group.

Other developments control the solution overgroup more closely.
Klyachko, Mikheenko, and Roman'kov obtain solutions within specified
classes of solvable groups with suitable torsion-free abelian factors
\cite{KlyachkoMikheenkoRomankov2024}. Ramirez-C\^ot\'e and Wise use
the Banach fixed-point theorem for groups embedded in Magnus-type
power-series groups; their nonsingular-system construction inverts the
determinant in the coefficient ring
\cite[Theorems~9--10]{RamirezCoteWise2025}. Such conclusions address
additional structure beyond existence in an arbitrary overgroup.

The direct antecedent of our proof is the spectral-phase and planar
argument in \emph{The Kervaire theorem for groups}
\cite[Lemma~2.1, Lemma~3.1, and Theorem~4.1]{OpenAIKervaire2026}.
Its fixed-plane incidence construction treats one unimodular relator.
The extension below uses products of unitary groups and rational
intersection classes to accommodate several relator types with
different multiplicities. The cohomological background is classical
\cite{Borel1953}; related geometry of unitary eigenvalue strata and
their intersection classes is developed by Nicolaescu
\cite{Nicolaescu2010}. In finite dimensions, spectral-phase
subadditivity also follows from Thompson's exponential formula
\cite{Thompson1986}. Antezana, Larotonda, and Varela prove an
approximate version, in operator norm, for embeddable $\mathrm{II}_1$ factors
\cite[Theorem~4.4]{AntezanaLarotondaVarela2012}. Our argument instead
uses the faithful trace of the group von Neumann algebra
\cite[Section~3.3]{Jones2009} and proves the required phase inequality
directly, so the coefficients need no unitary approximation hypothesis.


\subsection{Proof and technical contribution}

The proof passes from a possible kernel relation to a planar surface.
A finite product of conjugates of the relators is represented by disks
whose boundaries read the relators or their inverses. Rectangular bands
pair occurrences of the same variable with opposite signs. The remaining
boundary arcs carry only coefficients from $G$. The key assertion is a
boundary obstruction: on a connected planar surface of this form, if all
but one boundary words are trivial in $G$, then the last is trivial as
well. Removing the innermost components then proves injectivity.

Two independent estimates establish the obstruction. First, full row rank
forces equal numbers of positive and negative disks of each relator type.
Suppose these numbers are $k_i$ for the types present. A topological
argument supplies unitary matrices mixing the $k_i$ copies of each type
so that a finite-dimensional unitary matrix has at least $2\sum_i k_i$ fixed
directions. Second, a spectral-phase function measures the boundary words
through the left regular representation of $G$. Its subadditivity turns
the fixed-space estimate into an upper bound. A direct calculation around
the boundary circles attains that bound plus a nonnegative term measuring
the possibly nontrivial boundary word. Euler characteristic makes the
two constants agree, and faithfulness of the trace forces that word to
be the identity.

The spectral-phase and planar-surface strategy comes from the
one-relator unimodular theorem in~\cite{OpenAIKervaire2026}. We give
the required arguments in full. The principal extension is the
multi-block fixed-space theorem in Section~\ref{sec:fixed}. It concerns
block matrices built from independent unitary variables of possibly
different sizes, followed by arbitrary fixed unitary matrices. Full row
rank of the signed block-multiplicity matrix forces a simultaneous lower
bound on the dimensions of their fixed spaces. The proof pairs
fixed-subspace incidence cycles with the product of the unitary-variable
groups. Only the maximal-length exterior terms in cohomology contribute;
their coefficient is a product of nested nonzero minors. This makes
rational nonsingularity sufficient and gives an intersection principle
that is independent of its application to equations over groups.

Section~\ref{sec:phase} establishes the phase inequality in the finite
operator algebra of the left regular representation and computes the
phase of a weighted cyclic shift. Section~\ref{sec:planar} combines
these results to prove the boundary obstruction. Section~\ref{sec:kernel}
constructs the planar surfaces from an arbitrary kernel relation and
completes the proof. The unitary topology is used only for finite complex
matrices on the copy indices; the group coefficients are retained as
operators throughout. This separation allows the argument to apply to
arbitrary coefficient groups.
